\begin{remark}[What the corollary reaches, and the two facts that close the gap]\label{rem:corollary-reach}
``Each is an instance'' is the part that needed care, and an earlier version of this
sentence overstated it. The corollary is about three \emph{specific} languages and is
closed under neither concatenation with a literal nor union, while most refused corpus
patterns are of those shapes: \texttt{dm-.+} is not $\Sigma^{+}$, and
\texttt{5[0-9]+} is not $A^{+}$, because $A^{+}$ contains words of length one and that
language's shortest word has three characters. Two facts close the gap, and both are
already in this section rather than new.

The first is the uniqueness step inside Proposition~\ref{prop:decidable}'s proof: if
$L(r)=L(g)$ then $L(r)$ has a \emph{unique} shortest word. Most of the patterns the
corollary does not reach fail that immediately --- \texttt{dm-.+} has shortest words
\texttt{dm-} and nothing else of that length only if the suffix can be empty, which
\texttt{.+} forbids, so its shortest words are \texttt{dm-}$\sigma$ for every $\sigma$,
and there are $|\Sigma|$ of them. The second is the generalisation part~(3) is an
instance of, which is now Corollary~\ref{cor:absent-char}: if $L$ is infinite and some
$\sigma\in\Sigma$ occurs in no word of $L$, then $L$ is no union of finitely many globs.
That is what the prose below had always said, and the tenth pass is where it stopped
being an aside and became the numbered statement the $10$ unconditional occurrences
rest on.
\end{remark}
